\scanpage{116}
on a ici
\[
\rho=\frac{b^2}{c}-2,\qquad\overline\rho=\frac{\overline b^2}{\overline c}-2
\]
la condition $\rho=\overline\rho$ s'écrit
\[
\frac{b^2}{c}=\frac{\overline b^2}{\overline c}\quad,
\]
qui n'implique pas l'isom. des deux extensions ! P. ex. en car. $\ne2$, si on prend $b=\overline b=0$ la condition est vérifiée \deleted{($\rho=-2$} (il est vrai que l'on est dans le cas exceptionnel $\rho=\unsure{-|+}2$) mais l'on n'a d'isom entre les \deleted{\unsure{rev.|ext.}} $\unsure{\xi|\mathcal E|\varepsilon}_H$ et $\xi_{\overline H}$ que si $\overline c/c$ est un carré. \deleted{\unsure{et\ldots}} Etc\ldots

Il serait intéressant de voir si tout $\rho_0$ inv. peut être réalisé loc. \uuline{Zar} (et non seulement fpqc, ou fppf) Je soupçonne que oui, en partant d'un $H$, et en prenant $H'=\operatorname{int}(u)H$ pour $u\in A^{\unsure{\times|*}}$ convenable\ldots.

\uline{NB} Vérification immédiate dans le cas où $A$ est scindée\ldots
% archive sheet number: 115
